\scanpage{65}
% archive foliation: 67
[On laisse tomber dorénavant les $'$ des $X'$ !]

\deleted{Mais sur $X^*=X\setminus\bigcup D_i$ (si \unsure{chaque composante}
\unsure{connexe} \unsure{$X$ est}}

OPS. $X$ connexe, $I$ non vide --- alors
il y a sur $X^*=X\setminus\bigcup_i D_i^\circ$ (surface
conforme à bord obtenue dans la
découpe de $X$ suivant $\bigcup\delta D_i$) une
\deleted{\unsure{mét}} métrique \deleted{hyperbolique} canonique modulo
constante (qu'on peut d'ailleurs normaliser
par la condition que le volume soit $1$)
\deleted{\unsure{sauf dim $=$}}\inserted{et} canonique \inserted{sans cte} dans le cas
hyperbolique (i.e. quand on n'a
\[
g=0\qquad \underset{\substack{\Vert\\\text{card }I}}{\nu}\in\{1,2\}
\]
) --- dans le cas
\margin{il y a aussi le cas \deleted{\unsure{dans}} \unsure{orienté}
du plan projectif réel $(\simeq S^2/\pm1)$
\deleted{\unsure{privé} d'un disque}\inserted{``avec un trou''}}
hyperbolique le bord est géodésique,
donc dans tous cas le bord est muni
d'une structure torique canonique
provenant de la structure riemannienne.
Il est tentant de restreindre le \inserted{syst. des}
disques envisagés, en exigeant que
les deux structures toriques sur
chaque $\delta D_i$ (provenant soit de $D_i$, soit de $X^*$)
